\documentclass[11pt,a4paper]{article}
\usepackage[utf8]{inputenc}
\usepackage{amsmath,amssymb,amsthm}
\usepackage{listings}
\usepackage{geometry}
\usepackage{hyperref}
\geometry{margin=2.5cm}

\lstset{
  language=C++,
  basicstyle=\ttfamily\small,
  keywordstyle=\bfseries,
  breaklines=true,
  frame=single,
  numbers=left,
  numberstyle=\tiny,
  tabsize=4,
  showstringspaces=false
}

\title{IOI 2006 -- Forbidden Patterns (Writing)}
\author{Solution}
\date{}

\begin{document}
\maketitle

\section{Problem Statement}
Given a set of $K$ forbidden patterns (short strings) over an alphabet of size $\Sigma$, count the number of strings of length $N$ that do not contain any forbidden pattern as a substring. Output the count modulo a given prime.

\section{Solution}

\subsection{Aho-Corasick Automaton + DP}
We build an Aho-Corasick automaton from all forbidden patterns and then run a DP over the automaton states.

\begin{enumerate}
    \item \textbf{Build the automaton.} Insert all forbidden patterns into a trie, then compute failure links via BFS from the root. During BFS, propagate the ``forbidden'' flag: a state is forbidden if it is the endpoint of a pattern or its failure chain reaches such a state.
    \item \textbf{Complete the goto function.} For each state $s$ and character $c$, if $s$ has no child on $c$, set $\delta(s,c) = \delta(\text{fail}(s), c)$. This eliminates the need to follow failure links during the DP.
    \item \textbf{DP.} Let $dp[i][s]$ be the number of valid strings of length $i$ that leave the automaton in state $s$. The base case is $dp[0][\text{root}] = 1$. The transition is:
    \[
    dp[i+1][s'] \mathrel{+}= dp[i][s] \qquad \text{for each character } c \text{ with } \delta(s,c) = s',
    \]
    skipping transitions where $s$ or $s'$ is forbidden.
    \item \textbf{Answer.} $\displaystyle\sum_{s \text{ not forbidden}} dp[N][s]$.
\end{enumerate}

\subsection{Matrix Exponentiation (for large $N$)}
The DP transition can be expressed as a matrix-vector product. Let $M$ be the number of non-forbidden states and $A$ the $M \times M$ transition matrix. Then $dp[N] = A^N \cdot dp[0]$, computable in $O(M^3 \log N)$ via fast matrix exponentiation.

\subsection{Complexity}
\begin{itemize}
    \item \textbf{Basic DP:} $O(N \cdot M \cdot \Sigma)$ time, $O(M)$ space (rolling two DP layers).
    \item \textbf{Matrix exponentiation:} $O(M^3 \log N)$ time, $O(M^2)$ space.
\end{itemize}
Here $M = O(\sum |P_i|)$ is the number of automaton states.

\section{C++ Solution}

\begin{lstlisting}
#include <bits/stdc++.h>
using namespace std;

const int MOD = 1e9 + 7;

struct AhoCorasick {
    int ALPHA;
    struct Node {
        vector<int> ch;
        int fail;
        bool forbidden;
        Node(int alpha) : ch(alpha, -1), fail(0), forbidden(false) {}
    };
    vector<Node> trie;

    AhoCorasick(int alpha) : ALPHA(alpha) {
        trie.emplace_back(ALPHA); // root = node 0
    }

    void insert(const string& s) {
        int cur = 0;
        for (char c : s) {
            int ci = c - 'a';
            if (trie[cur].ch[ci] == -1) {
                trie[cur].ch[ci] = trie.size();
                trie.emplace_back(ALPHA);
            }
            cur = trie[cur].ch[ci];
        }
        trie[cur].forbidden = true;
    }

    void build() {
        queue<int> q;
        for (int c = 0; c < ALPHA; c++) {
            if (trie[0].ch[c] == -1)
                trie[0].ch[c] = 0;
            else {
                trie[trie[0].ch[c]].fail = 0;
                q.push(trie[0].ch[c]);
            }
        }
        while (!q.empty()) {
            int u = q.front(); q.pop();
            trie[u].forbidden |= trie[trie[u].fail].forbidden;
            for (int c = 0; c < ALPHA; c++) {
                if (trie[u].ch[c] == -1)
                    trie[u].ch[c] = trie[trie[u].fail].ch[c];
                else {
                    trie[trie[u].ch[c]].fail = trie[trie[u].fail].ch[c];
                    q.push(trie[u].ch[c]);
                }
            }
        }
    }

    int size() const { return trie.size(); }
};

int main(){
    ios::sync_with_stdio(false);
    cin.tie(nullptr);

    int N, K, ALPHA;
    cin >> N >> K >> ALPHA;

    AhoCorasick ac(ALPHA);
    for (int i = 0; i < K; i++) {
        string p;
        cin >> p;
        ac.insert(p);
    }
    ac.build();

    int M = ac.size();

    // DP with rolling array
    vector<long long> dp(M, 0);
    dp[0] = 1;

    for (int i = 0; i < N; i++) {
        vector<long long> ndp(M, 0);
        for (int s = 0; s < M; s++) {
            if (dp[s] == 0 || ac.trie[s].forbidden) continue;
            for (int c = 0; c < ALPHA; c++) {
                int ns = ac.trie[s].ch[c];
                if (!ac.trie[ns].forbidden)
                    ndp[ns] = (ndp[ns] + dp[s]) % MOD;
            }
        }
        dp = ndp;
    }

    long long ans = 0;
    for (int s = 0; s < M; s++)
        if (!ac.trie[s].forbidden)
            ans = (ans + dp[s]) % MOD;

    cout << ans << "\n";
    return 0;
}
\end{lstlisting}

\section{Notes}
The Aho-Corasick automaton captures exactly the suffix information needed to detect forbidden patterns. By completing the goto function during construction, the DP transition becomes a simple table lookup. The forbidden-flag propagation through failure links ensures that states reachable via a suffix match are correctly marked.

For the IOI 2006 constraints, the basic $O(NM\Sigma)$ DP suffices. If $N$ is extremely large (e.g., $N \le 10^{18}$), use matrix exponentiation on the $M \times M$ transition matrix.

\end{document}
